\subsection{Integral orthogonal complements and the actual compact cross boundary}

\begin{statement}{D.63}{Proved}{An integral compact correction lifts fifty-two more arithmetic invariant coordinates}
Keep the actual sources and maps of D.58--D.62:
\[
 X=\widehat B/(R_{50}\cap L_2),\qquad
 Y=X/\mathcal K,\qquad X_{62}=X/\mathcal J_{62}.
 \tag{D.63.1}
\]
Let \(P=U/\mathcal R\) be integral compact homology,
with its actual perfect alternating intersection form.
On each of the four small-character orbits of D.60,
write its already constructed compact decomposition
as \(P_a=E_a\oplus M_a\). Here \(E_a\) is the sum
of the four extra projective summands, of fibre rank
32 over the unramified quadratic coefficient ring,
and \(M_a\) has fibre rank 16.

The restriction of the intersection form to every
\(M_a\) is perfect. There are exact integral
\(G\)-equivariant graph isomorphisms
\[
 \psi_a:E_a\xrightarrow{\ \sim\ }\widetilde E_a=M_a^\perp,
 \qquad P_a=\widetilde E_a\mathbin{\perp}M_a.
 \tag{D.63.2}
\]
They retain the kernel of each old compact projector.
They have affine lifts inside the actual kernel
\(\mathcal L=\ker(U^-\to U_*^-)\), and every old
central projector kernel used in \(\mathcal J_{62}\)
is unchanged. The proof supplies the actual graph,
affine inclusion and retraction maps; it does not
identify two differently marked embeddings without
these maps.

Let \(\widetilde e_{a,i}\) be the resulting affine
projectors onto the four summands of
\(\widetilde E_a\), and let \(m_a\) be the compatible
affine projector onto \(M_a\). Both kill all cusp
homology. With the cross-projector convention of D.60,
put
\[
 \begin{aligned}
 \Pi_{63}&=e_\Delta\sum_{a=1}^4\sum_{i=1}^4
                   (\overline{\widetilde e}_{a,i}\boxtimes\bar m_a),\\
 \mathcal C_{63}&=\Pi_{63}Z,\qquad
 \mathcal J_{63}=(1-\Pi_{63})\mathcal J_{62}.
 \end{aligned}
 \tag{D.63.3}
\]
These are actual equivariant central submodules, and
\[
 \begin{gathered}
 \mathcal C_{63}\subset\mathcal J_{62},\qquad
 \mathcal J_{62}=\mathcal C_{63}\oplus\mathcal J_{63},\\
 \dim\mathcal C_{63}=40960,\qquad
 \dim(\mathcal C_{63})_G=52,\\
 \nu:(\mathcal C_{63})_G\xrightarrow{\ \sim\ }\mathcal C_{63}^G,
 \qquad \bigwedge^2\pi\text{ is injective on }\mathcal C_{63}.
 \end{gathered}
 \tag{D.63.4}
\]
The exact orbit contributions are
\[
 \begin{array}{c|r|r|c}
 \ell&\dim\mathcal C_{63,\ell}&\dim(\mathcal C_{63,\ell})_G
     &\text{four individual norm ranks}\\\hline
 (0,0,0,1)&8192&12&(4,4,4,0)\\
 (0,0,1,1)&12288&16&(4,4,4,4)\\
 (0,0,1,2)&12288&16&(4,4,4,4)\\
 (0,1,0,0)&8192&8&(2,2,2,2)
 \end{array}
 \tag{D.63.5}
\]

There is an actual continuous peripheral automorphism of
\[
                         X_{63}=X/\mathcal J_{63}.
 \tag{D.63.6}
\]
It has multiplier one, fixes the original cusp \(t\),
is identity on \(G\), and retains the complete D.57
action with its literal compatible conjugators.
Its full \(\mathcal W_h\) action, normalized parameter
and all established metabelian and CM comparisons
through that map are unchanged.

Here is the exact change of representative on \(Y\).
Write \(T=T_{61}\) for the integral homology action
of D.62 and \(e=\sum_{a,i}\widetilde e_{a,i}\).
The new integral homology action is
\[
                    T_{63}=(1-e)T(1-e)+e.
 \tag{D.63.7}
\]
It is an automorphism and is identity on all
\(\widetilde E_a\). There is a continuous derivation
\(\eta:Y\to eU\), induced from \(\mathcal J_\infty\),
whose restriction on \(U=\ker(\mathcal J_\infty\to G)\)
is \(e\), with \(\eta(t)=0\) and peripheral coboundaries.
Define
\[
 \begin{aligned}
 A_0&=-(1-e)T|_{eU}:eU\longrightarrow\mathcal L,\\
 \beta_Y(g)&=(A_0\eta(g))\alpha_{62,Y}(g).
 \end{aligned}
 \tag{D.63.8}
\]
The coefficient in the second line is the actual
element of the abelian kernel \(\mathcal L\subset Y\).
After the self-block corrections of D.62, the resulting
\(Y\)-action is
\[
                  \operatorname{Int}(b_{t,Y})\beta_Y,
 \qquad
 b_t=\prod_j\prod_{k=0}^{35}t^ku_jt^{-k},\quad u_j\in S_{57}.
 \tag{D.63.9}
\]
The vectors of the \(u_j\) are the transported D.62
correction vectors selected by the actual boundary.
The element \(b_{t,Y}\) centralizes \(t\), has zero
image in \(X_*\) and \(\mathcal W_h\), and acts
trivially on integral homology. Thus the projection
to \(X_{62}\) is an actual newly constructed action;
(D.63.8)--(D.63.9) specify its comparison with the
chosen D.62 representative.

All source kernels are retained explicitly:
\[
 \begin{aligned}
 \ker(X_{63}\to X_{62})&=\mathcal C_{63},\\
 \ker(X_{63}\to Y)&=(1-e_\Delta)\mathcal K
       \oplus\mathcal C_{59}\oplus\mathcal C_{60}
       \oplus\mathcal C_{61}\oplus\mathcal C_{62}\oplus\mathcal C_{63},\\
 \ker(X\to X_{63})&=\mathcal J_{63},\\
 \dim\mathcal J_{63}&=48w-2560-5308416\kappa-81920.
 \end{aligned}
 \tag{D.63.10}
\]
The specified small-character residual has dimension
284160 and coinvariant dimension 482. Its compact
norm has rank six and kernel dimension 476. With
\[
 S_{63}=\{b\in R_{50}:
       \operatorname{im}_{\mathcal T_2}(b)\in\mathcal J_{63}\},
 \qquad S_{63}/(R_{50}\cap L_2)=\mathcal J_{63},
 \tag{D.63.11}
\]
this is a nonzero remaining actual source kernel.
Its boundary and the full profinite realization
remain unresolved. The manuscript passages are
Part D(e), section 51.16, folio 734 verso, scan 174,
and section 51.15, folios 725--727, scans 156--160.
\end{statement}
\begin{proof}
\emph{The compact markings and the projective types.}
All maps start from the marked arithmetic quotient
\(G\to G^\dagger=(\mathbb Z/3)^4\rtimes Q\to G_*\)
of D.60.2. The same forty unordered characters are
used, with their original pair-coset representatives.
On one character fibre, the actual affine module has
rank 96 over \(\mathbb F_4\), its cusp submodule has
rank 48, its compact quotient has rank 48, and its
small compact target has rank four.

The program \path{d_compact_splitting_maps.py} checks
the actual compact quotient map, the small map and
their square on every matrix column. For the old
orbits the retained compact decomposition is exactly
\([E_0,E_1,P_0]\), with ranks \((16,16,16)\).
The last block is \(M_a\); the first two split into
the four rank-eight extra summands. The program
checks these positions against both recorded actual
affine embeddings and the compact target map.
For the new orbits it uses the D.61 rank-32
projective embedding and its recorded rank-16
complement, and checks all stabilizer intertwiners.

The full unordered Sylow orders, real dimensions
and norm ranks are
\[
 \begin{array}{c|r|r|r|r|r|r|r}
 \text{orbit size}&|P_2|&\dim E&\operatorname{rank}N_{P_2}|E
   &\dim M&\operatorname{rank}N_{P_2}|M
   &\dim\mathcal R&\operatorname{rank}N_{P_2}|\mathcal R\\\hline
 8&8&64&8&32&4&96&12\\
 12&16&64&4&32&0&96&6
 \end{array}
 \tag{D.63.12}
\]
This uses the semilinear exchange element as well as
the oriented stabilizer. To justify the projectivity
test, choose vectors with independent norm images.
They define a map from the corresponding number of
regular modules. A nonzero kernel would contain a
nonzero invariant vector, since every nonzero module
of a finite two-group in characteristic two has one.
The invariant vectors of a sum of regular modules
are precisely its independent norm vectors, so that
kernel is zero. Equality of dimensions makes the
map an isomorphism. Restriction to a Sylow subgroup
detects projectivity: induction followed by the
coset trace splits the original module because the
index is odd. This proves projectivity of all the
cusp modules and all extra summands in (D.63.12).
Induction over the actual pair cosets and odd
averaging along the character subgroup and
\(G\to G^\dagger\) give their projective modules
over the full actual group \(G\).

In the new orbits the norm on the oriented
Sylow \(D_8\) of the complement is already zero,
as checked in D.61. A nonzero projective summand
would restrict to a nonzero free module with
nonzero norm. Hence these complements have no
projective direct summand.

In the old orbits both pieces are projective, but
their indecomposable projective types are disjoint.
The normal \(V_4\)-coinvariant top is a module for
the odd complement \(C_3^2\). The extra and cover
weight supports, in the original two-generator
marking, are
\[
 \begin{array}{c|c|c}
 \ell&\text{extra support}&\text{cover support}\\\hline
 (0,0,0,1)&\{(0,1),(0,2),(1,0),(2,0)\}
                       &\{(1,2),(2,1)\}\\
 (0,1,0,0)&\{(0,1),(0,2),(1,1),(2,2)\}
                       &\{(0,0),(1,2),(2,1)\}.
 \end{array}
 \tag{D.63.13}
\]
The normal \(V_4\) augmentation ideal is nilpotent,
and its quotient algebra is the split semisimple
algebra of \(C_3^2\). Thus these are exactly the
simple types of the projective tops. Their duals
negate the weights: the \(V_4\) socle is its norm,
which the odd complement preserves. Each extra
support is closed under negation and is disjoint
from its cover support. Consequently the old
\(E_a^*\) and \(M_a\) have no common projective
direct summand. The displayed weights and their
multiplicities are recomputed from all actual
matrices in the certificate.

\emph{The perfect restricted form and the exact graph.}
Work first on the actual integral compact lattices.
The character-orbit summands are mutually orthogonal:
invariance under the odd character group pairs a
character only with its inverse. Since each entire
unordered orbit is retained, the intersection form
on \(P_a\) is perfect.

Write its isomorphism \(P_a\to P_a^*\) and inverse
in blocks for \(P_a=E_a\oplus M_a\). If the lower
diagonal blocks are \(D_a:M_a\to M_a^*\) and
\(D_a':M_a^*\to M_a\), multiplication gives
\[
                     D_a'D_a=1_{M_a}-C_a'B_a.
 \tag{D.63.14}
\]
The last endomorphism factors through \(E_a^*\).
Modulo two, an invertible Fitting summand of it
would split off from both \(M_a\) and \(E_a^*\):
compose the two factor maps with the inverse on
that summand to obtain the retraction and inclusion.
For the new orbits this contradicts the absence
of projective summands in \(M_a\). For the old
orbits it contradicts (D.63.13), after restriction
to the oriented stabilizer. The Fitting summand
is therefore zero, the factor endomorphism is
nilpotent, and (D.63.14) is invertible modulo two.
It follows that \(D_a\) is invertible integrally.

Let \(C_a:E_a\to M_a^*\) be the lower off-diagonal
block of the intersection form. The actual map in
(D.63.2) is
\[
       \psi_a(v)=v-\iota_{M_a}D_a^{-1}C_a(v).
 \tag{D.63.15}
\]
It has the same \(E_a\)-coordinate as \(v\), pairs
to zero with all \(M_a\), and is an isomorphism
onto \(M_a^\perp\). In particular the new compact
projector is \(\psi_a e_{a,c}\), where \(e_{a,c}\)
is the old compact projector. Its kernel is still
\(M_a\). Each of the four individual projectors
is transported by this same graph map.

Let \(\theta_P:P\to P_*\) be the actual compact
map. The perfect forms define the equivariant
adjoint \(\theta_P^\dagger:P_*\to P\), by
\[
       \langle\theta_Pv,w\rangle_*
                  =\langle v,\theta_P^\dagger w\rangle.
 \tag{D.63.16}
\]
For each of the four fibres the certificate forms
every coefficient equation
\(E_gH=H(P_*)_g\) for the actual oriented generators.
Their rank is 128 on the 128-dimensional
\(\mathbb F_4\) space of maps from the rank-four
target to the rank-32 extra part. Thus
\(\operatorname{Hom}_G(\bar P_*,\bar E_a)=0\).
An integral equivariant map reducing to zero is
twice another such map, since the target lattice
is free. Iteration and separatedness give
\(\operatorname{Hom}_G(P_*,E_a)=0\).
Consequently \(e_{a,c}\theta_P^\dagger=0\), so its
image in this orbit lies in \(M_a\). Equations
(D.63.15)--(D.63.16) and perfection on \(P_*\)
give the exact conclusion
\[
                         \theta_P\widetilde E_a=0.
 \tag{D.63.17}
\]

\emph{Affine inclusions inside the retained source kernel.}
The integral cusp map is onto. Indeed the actual
orders of the marked cusp in \(G\) and \(G_*\)
are 36 and 12. A source boundary loop
\(g t^{36}g^{-1}\) maps to the cube of its target
boundary loop. Every target cusp has a source
preimage; multiplication by three is a unit over
\(\mathbb Z_2\), and the single sum relation maps
to the corresponding sum relation. These literal
boundary maps prove surjectivity, also after the
odd character projectors. The finite certificate
checks its rank four on every target cusp fibre.

The diagram with exact rows
\(0\to\mathcal R\to U\to P\to0\) and its small
counterpart therefore gives an onto map
\[
 \ker\theta_U\longrightarrow\ker\theta_P.
 \tag{D.63.18}
\]
Its restriction to the active part is the same
map with source \(\mathcal L\). Projectivity of
\(\widetilde E_a\), which is integrally isomorphic
to \(E_a\), lifts its inclusion in (D.63.17)
through (D.63.18). Denote the resulting inclusion
by \(j_a:\widetilde E_a\hookrightarrow\mathcal L\).
Compactification is exactly the prescribed inclusion,
so this is a split injection of free lattices.
This proves the affine small-map zero without
postulating a stronger affine Hom vanishing.

Write \(U_\Sigma,\mathcal R_\Sigma,P_\Sigma\) for
the sums of these four character components.
On this specified part the cusp lattice is integrally
projective. Its projective reduction in (D.63.12)
lifts to an idempotent projective lattice; projectivity
lifts an isomorphism modulo two onto the actual
cusp lattice. The resulting map is an isomorphism
by its invertible reduction. The dual cusp lattice
is also projective, since the dual of a regular
group-ring lattice is regular. Dualizing the affine
exact sequence therefore splits it, and dualizing
back supplies a retraction
\(r:U_\Sigma\to\mathcal R_\Sigma\).
This argument requires no projectivity of \(P\).

Replace this retraction by
\[
                  r'=r-\sum_a rj_a\widetilde e_{a,c}\pi.
 \tag{D.63.19}
\]
It is still identity on \(\mathcal R_\Sigma\) and vanishes
on all \(j_a\widetilde E_a\). Thus \(1-r'\)
induces a compact section
\(j_P:P_\Sigma\to U_\Sigma\) agreeing
with the chosen \(j_a\) on these summands. Set
\[
 \widetilde e_{a,i}=j_a\psi_a e_{a,i,c}\pi,
 \qquad m_a=j_P(1-\widetilde e_{a,c})\pi
       \quad\text{on the }a\text{th orbit}.
 \tag{D.63.20}
\]
Extend the projectors by zero on all other character
components. These are mutually orthogonal equivariant affine
idempotents; they kill cusps and inject under
compactification. Their reductions have precisely
the old module matrices in the transported bases.

There is an explicit comparison with the old affine
marking. Let \(e^{\rm old}\) be the sum of its
sixteen old projectors, and use the old compact
inclusion to interpret \(\psi\pi e^{\rm old}\).
Then
\[
          \Psi=1-e^{\rm old}+j\psi\pi e^{\rm old}:U\to U
 \tag{D.63.21}
\]
is an integral equivariant automorphism, with
\((\Psi-1)^2=0\), \(\theta_U\Psi=\theta_U\), and
\(\Psi|_{\mathcal R}=1\). Its new projectors satisfy
\(\widetilde e_{a,i}=\Psi e_{a,i}^{\rm old}\),
and the old projectors satisfy
\(e_{a,i}^{\rm old}\Psi=e_{a,i}^{\rm old}\).
The exterior map \(\bigwedge^2\Psi\) carries every
old selected self or cross image to its new image
and has the same kernel on that projector.
It also satisfies
\(\theta_Z\bigwedge^2\bar\Psi=\theta_Z\), so it
preserves the actual comparison kernel \(\mathcal K\).
It fixes the other character orbits. In particular
the literal submodule \(\mathcal J_{62}\) is unchanged.

\emph{The integral homology action and its actual peripheral realization.}
D.60--D.62 give
\(e_{a,i}^{\rm old}T=e_{a,i}^{\rm old}\).
Hence \(eT=e\) for the new total projector.
The kernel \(\ker e\) is invariant under \(T\),
and its restriction is an automorphism, since
\(T\) is an automorphism and induces identity on
the quotient \(eU\). Relative to
\(U=eU\oplus\ker e\), its matrix is
\(\left(\begin{smallmatrix}1&0\\ C&D\end{smallmatrix}\right)\).
Formula (D.63.7) replaces this by
\(\left(\begin{smallmatrix}1&0\\0&D\end{smallmatrix}\right)\).
Thus \(T_{63}\) is invertible, fixes all cusps,
retains \(\theta_UT_{63}=T_*\theta_U\), and its
difference from \(T\) is \(A_0e\).
The image of \(A_0\) is in \(\mathcal L\) because
\(\theta_U eU=0\) and the old small square commutes.

The projective module \(eU\) has zero positive
cohomology over the finite group \(G\).
Inflation--restriction therefore lifts the actual
equivariant map \(e:U\to eU\) to a continuous
derivation of \(\widehat B\). Its value on
\(t^{36}\) is zero, since \(e\) kills cusp homology.
On the projective cyclic module, the exact norm
sequence gives
\(\ker N_{36}(t)=(1-t)eU\); subtract the corresponding
inner derivation to make its value on \(t\) zero.
The same projective norm sequence for \(s\) and
\(r\) gives actual peripheral values
\[
 \eta(s)=(1-s)d_s,\qquad
 \eta(r)=(1-r)d_r,\qquad \eta(t)=0,
             \qquad d_s,d_r\in eU.
 \tag{D.63.22}
\]
All these assertions hold over \(\mathbb Z_2\):
the regular-module cyclic identities split, and
direct summands preserve their exactness.
The derivation kills the kernel of the full
homology source \(\mathcal J_\infty\) of D.50,
so it factors through \(Y\).

Composing with the actual map \(A_0:eU\to\mathcal L\)
gives a derivation with values in the abelian
kernel of \(Y\to X_*\). Since \(\alpha_{62,Y}\)
is identity on \(G\), the multiplication formula
in (D.63.8) is a homomorphism. Its peripheral
conjugators change by \(A_0d_s,A_0d_r\), which
have zero images in \(X_*\) and \(\mathcal W_h\).
It fixes \(t\) and induces the invertible matrix
\(T_{63}\). Its maps on the quotient and on the
kernel in the D.58 fibre product are automorphisms,
so it is an actual automorphism of \(Y\).
Each new conjugating element has an actual lift
in \(S_{57}\) by D.58.2. Choosing these lifts gives
the literal source conjugators used below.

\emph{The cross-module norms and the actual boundary zero.}
In a transported fibre, a cross block has coordinates
\(\operatorname{Mat}_{8\times16}(\mathbb F_4)\),
of dimension 256 over \(\mathbb F_2\). An oriented
stabilizer element acts by
\(H\mapsto E_gH\overline{M_g}^{\,t}\); a character
exchange applies coefficient conjugation as well.
The program \path{d_all_compact_cross_norm.py}
forms every such column, all relations \((g-1)H\),
and their complete group norm. It verifies that
the norm kills every relation and that every norm
column is invariant. The quotient dimension equals
the norm rank in each of all sixteen blocks,
with the values in (D.63.5). Induction over the
actual pair cosets multiplies dimensions by eight
or twelve and leaves coinvariants and norm ranks
as displayed. This proves (D.63.4).

The compact action of \(T_{63}\) is block diagonal
on the actual orthogonal decomposition
\(\widetilde E_a\perp M_a\), and is identity on
the first part. The compact orientation bivector
has zero cross projection there. Its image under
\(\bigwedge^2\overline T_{63}\) has the same zero
projection. For a preliminary source lift with
\(\varphi(t)=zt\), D.55's full, unquotiented
orientation identity therefore gives
\[
             (\bigwedge^2\pi)\nu([\Pi_{63}z])=0.
 \tag{D.63.23}
\]
Both maps on this selected module are injective
by (D.63.4), so the actual cross boundary class is
zero. This evaluates the boundary through its
actual orientation map, rather than transferring
a conclusion from the separate order-1080 example.

The chosen affine complement can have cusp terms
under \(T_{63}\). This does not affect descent:
its projected action satisfies
\[
 \widetilde e_{a,i}T_{63}=\widetilde e_{a,i},\qquad
 m_aT_{63}=D_{a,c}m_a,
 \tag{D.63.24}
\]
where \(D_{a,c}\) is the compact complement action
expressed on the chosen affine section. The second
formula follows by compactification; both sides
kill cusps and all other components. Thus the
kernel of each cross projector is invariant,
with invertible induced map on its image.

\emph{Retaining the earlier coordinates and correcting the source.}
The graph comparison (D.63.21) preserves every old
central projector kernel. Moreover the sixteen
new projectors have fixed left rows for \(T_{63}\).
The D.60 and D.61 selected self and cross norms
therefore still force zero boundary on their
injective-norm components. The D.59 selected
character orbits are unchanged.

For the remaining twenty-six self coordinates,
transport the exact D.62 maps by \(\Psi\):
\[
 [v]\longmapsto[e_\Delta(v\wedge sv)]_G,
 \qquad
 [v]\longmapsto[e_\Delta(v\wedge N_{36}(t)v)]_G.
 \tag{D.63.25}
\]
The homology vectors now lie in the actual
\(\widetilde E_{a,i}\subset\mathcal L\); their
matrices, induced coset bases and all twenty-six
independent columns are unchanged through the
proved module isomorphism. D.58.2 supplies their
actual preimages \(u_j\in S_{57}\).
The source involution words
\(u\widetilde s(u)\) and cusp words
\(\prod_{k=0}^{35}t^kut^{-k}\) are exactly those
of D.62, with the new source conjugators.
In particular the wrap term is still
\(t^{36}ut^{-36}\), not \(u\).
Their errors lie in the selected self block plus
terms with a cusp factor. Hence their projection
to \(\mathcal C_{63}\) is zero. They cancel the
whole self norm kernel, change the \(Y\)-action
by precisely (D.63.9), and retain its homology.

Every now retained central boundary class is zero.
The exact augmentation map
\((b,c)\mapsto(1-s)b+s(1-r)c\) realizes its remaining
representative as actual central changes of the
two conjugators. These corrections preserve \(Y\)
and the complete \(X_*\) conjugators. The final
cusp error lies in \(\mathcal J_{63}\).
The defining class-two laws and the projector
kernel identities give descent through that
submodule. The invertible Frattini map gives an
automorphism of the finite class-two component;
its common quotient with the full homology and
\(\mathcal W_h\) sources remains equivariant.
Their actual marked fibre product is exactly
\(X/\mathcal J_{63}\). This proves (D.63.6).

\emph{The actual remaining kernel and the computation record.}
The image of every cross projector is in
\(\mathcal J_{62}\): its first factor has zero
small image, its second is killed by every old
extra projector, and its character support is
disjoint from the D.59 support. Thus (D.63.3)
is a direct sum decomposition inside the literal
old source kernel. Subtracting its dimensions
gives
\[
            325120-40960=284160,\qquad
                         534-52=482.
 \tag{D.63.26}
\]
The original compact norm rank in D.60 was 138.
D.60, D.61 and D.62 removed ranks 28, 50 and 2,
respectively. The present compact-injective norm
image removes rank 52, leaving
\[
                138-28-50-2-52=6,
                   \qquad 482-6=476.
 \tag{D.63.27}
\]
The compact projectors make these images direct:
they kill the complementary residual and are
identity on their own compact norm images.
Consequently these are ranks of the actual
remaining comparison maps, not dimension estimates.
Taking a nonzero residual vector and its source
preimage under \(R_{50}\to Z\) gives (D.63.11)
and proves that the remaining kernel is nonzero.

The complete programs also retain a separate new-orbit
cross calculation and the cusp/complement module
calculation, with every original output. Two abandoned
old-orbit extractions are preserved in the trial record:
one assumed zero cross coinvariants, and another used
the wrong position of \(P_0\) in the recorded basis.
Neither is used in this proof. The final verifier
checks that position against the actual target and
affine maps before extracting any action. Its first
corrected mathematical run had a wall-clock variable
collision; that complete run is preserved, and the
fresh timing-corrected run has identical mathematical
fields. The separate splitting verifier's initial
missing input-key error is also retained. The final
four programs and all their complete outputs are
replayed independently from copied immutable inputs.
\end{proof}
